%% The Elicit manual design, on a few pages — every element elicit-brand.sty provides, in the
%% order a manual uses them. Typeset it after any change to the .sty and look at it before
%% rebuilding the four manuals: docs/manual/build-manual.sh --job specimen.
%%
%% It is the design's reference, not a manual: nothing in it describes the platform.

\documentclass[11pt, a4paper]{article}

\usepackage[a4paper, margin = 25mm, bottom = 28mm]{geometry}

\newcommand{\elicitProduct}{Elicit Design Specimen}
\usepackage{elicit-brand}
\usepackage[hidelinks]{hyperref}
\hypersetup{pdftitle = {Elicit — Manual design specimen}, pdfauthor = {Elicit Software LLC}}
\elicitAccent{elicitVerdigris}

\begin{document}

\elicitTitlePage{Manual Design Specimen}{Every element of the Elicit manuals}%
  {Title page \textperiodcentered\ parts \textperiodcentered\ chapters \textperiodcentered\
   call-outs \textperiodcentered\ blocks \textperiodcentered\ tables \textperiodcentered\
   diagrams \textperiodcentered\ screenshots}%
  {This specimen is the reference for elicit-brand.sty, the design the installation manual,
   the administrator's manual and Author's two manuals are typeset in.\\[0.5em]
   \textcopyright\ 2025--2026 Elicit Software LLC \textperiodcentered\ PolyForm Noncommercial License 1.0.0}

\pagenumbering{roman}
\begin{elicitnote}[Who this specimen is for]
Whoever changes \file{docs/manual/elicit-brand.sty}. Each page below uses the elements one
manual chapter would, so a change that spoils one of them shows here first.
\end{elicitnote}
\vspace{1.2em}
\elicitContents

\clearpage
\pagenumbering{arabic}

\elicitpart{Getting oriented}{What an Elicit site is made of, how its pieces depend on each
other, and what has to exist before you install any of them.}

\section{What Elicit is made of}
\label{sec:topology}

Elicit is a modular survey platform: a handful of applications that share one database, one
identity provider and one brand directory. A \term{site} is one installation of them. This
chapter names the pieces and how they connect, so that you can decide which of them your
site runs before you install anything. On screen you open \ui{Admin} \textgreater\
\ui{Apply Survey Definition}, and a long label such as \ui{Send a test email to this address
now} still breaks across a line like any other phrase in a paragraph.

\begin{elicitdiagram}{A site. Respondents reach Survey and administrators reach Admin. The
optional report services are called by Survey and Admin, never by a browser.}
  \node[elicit person, minimum width = 2.6cm] (resp) at (0, 0) {Respondent};
  \node[elicit person, minimum width = 2.6cm] (adm) at (6.4, 0) {Administrator};
  \node[elicit app, minimum width = 2.6cm] (survey) at (0, -1.8) {Survey};
  \node[elicit app, minimum width = 2.6cm] (admin) at (6.4, -1.8) {Admin};
  \node[elicit service, minimum width = 2.3cm] (idp) at (10.0, -1.25) {Identity provider};
  \node[elicit service, minimum width = 2.3cm] (smtp) at (10.0, -2.35) {SMTP relay};
  \node[elicit database, minimum width = 4.4cm] (db) at (3.2, -3.6)
    {PostgreSQL\\[-1pt]{\scriptsize the \val{survey} schema and the reporting schemas}};
  \node[elicit app, minimum width = 2.6cm] (fhhs) at (0, -5.6)
    {FHHS\\[-1pt]{\scriptsize\mdseries report service}};
  \node[elicit app, minimum width = 2.6cm] (ped) at (3.6, -5.6)
    {Pedigree\\[-1pt]{\scriptsize\mdseries draws pedigrees}};
  \begin{scope}[on background layer]
    \node[elicit group, fit = (fhhs) (ped), inner ysep = 12pt] (opt) {};
  \end{scope}
  \node[elicit group label, anchor = south west] at (opt.north west) {optional};
  \draw[elicit flow] (resp) -- node[elicit label] {access code} (survey);
  \draw[elicit flow] (adm) -- node[elicit label] {signs in} (admin);
  \draw[elicit flow] (survey) -- (db);
  \draw[elicit flow] (admin) -- (db);
  \draw[elicit flow] (admin.east) -- (idp.west);
  \draw[elicit flow] (admin.east) -- (smtp.west);
  \draw[elicit flow] ([xshift = -0.9cm]survey.south) -- node[elicit label, pos = 0.38]
    {reports} ([xshift = -0.9cm]fhhs.north);
  \draw[elicit flow] (fhhs) -- (db);
  \draw[elicit flow] (fhhs) -- (ped);
\end{elicitdiagram}

\subsection{The three call-outs}

\begin{elicitnote}[One address for the identity provider]
The issuer address you configure must resolve to the same provider \emph{from the browser}
and \emph{from inside the container}. The bundled example uses \val{host.docker.internal}.
\end{elicitnote}

\begin{elicitcaution}
A module that starts without the service token is not ready, and answers 503 on its
protected paths until it is given one.
\end{elicitcaution}

\begin{elicitsecurity}
The compose file gives every application the development token
\val{elicit-dev-service-token-not-for-a-real-site}. A real site generates its own with
\val{openssl rand -hex 32}.
\end{elicitsecurity}

\subsection{Commands, SQL and files}

Generate a token, then check that a module answers:

\begin{elicitcmd}
openssl rand -hex 32
curl -s http://localhost:8080/q/health/ready | jq '.status'
\end{elicitcmd}

\begin{elicitsql}
ALTER ROLE elicit_owner NOSUPERUSER;
GRANT CREATE ON DATABASE survey TO elicit_owner;
\end{elicitsql}

\begin{elicitfile}[fontsize = \footnotesize]
survey:
  environment:
    ELICIT_SERVICE_TOKEN: "elicit-dev-service-token-not-for-a-real-site"
    ELICIT_SERVICE_TRUSTED_ORIGINS: "http://host.docker.internal:8082"
\end{elicitfile}

\section{Installing a survey definition}
\label{sec:apply}

A procedure's steps are numbered in discs, and a figure the procedure points into carries
the same discs where the controls are.

\begin{enumerate}
  \item Sign in to Admin and open \ui{Apply Survey Definition} \elicitstep{1}.
  \item Read the setup banners \elicitstep{2}; each one links to the page that clears it.
  \item Choose \ui{Upload} \elicitstep{3} and pick the \file{.elicit} file Author exported.
\end{enumerate}

\screenshot[\elicitmark{1}{0.035}{0.45}\elicitmark{2}{0.985}{0.205}\elicitmark{3}{0.205}{0.555}]%
  {09-apply-definition.png}{Apply Survey Definition, before a file is chosen.}

\elicitpart{Reference}{Every setting Survey and Admin define, with its default and whether a
site has to change it.}

\section{Configuration reference}
\label{sec:reference}

The \emph{Change} column says what a site does with a setting: \chgset\ it, \chgmay\ change
it, or \chgkeep\ it.

\begin{settings}{tab:survey}{Settings Survey defines}
\opt{quarkus.datasource.jdbc.url} & \val{jdbc:postgresql://db:5432/survey} & \chgset & The database Survey connects to. A site points it at its own cluster. \\
\opt{quarkus.datasource.username} & \val{survey_user} & \chgmay & The role Survey connects as; it owns nothing and reads and writes the \val{survey} schema. \\
\opt{elicit.service.token} & \val{(none)} & \chgset & The bearer token the modules send each other. At least 32 characters; generate one with \val{openssl rand -hex 32}. \\
\opt{quarkus.flyway.migrate-at-start} & \val{true} & \chgkeep & Survey creates and upgrades its schema when it starts. \\
\opt{quarkus.datasource.jdbc.url} & \val{jdbc:postgresql://db:5432/survey} & \chgset & The database Survey connects to. A site points it at its own cluster. \\
\opt{quarkus.datasource.username} & \val{survey_user} & \chgmay & The role Survey connects as; it owns nothing and reads and writes the \val{survey} schema. \\
\opt{elicit.service.token} & \val{(none)} & \chgset & The bearer token the modules send each other. At least 32 characters; generate one with \val{openssl rand -hex 32}. \\
\opt{quarkus.flyway.migrate-at-start} & \val{true} & \chgkeep & Survey creates and upgrades its schema when it starts. \\
\opt{quarkus.datasource.jdbc.url} & \val{jdbc:postgresql://db:5432/survey} & \chgset & The database Survey connects to. A site points it at its own cluster. \\
\opt{quarkus.datasource.username} & \val{survey_user} & \chgmay & The role Survey connects as; it owns nothing and reads and writes the \val{survey} schema. \\
\opt{elicit.service.token} & \val{(none)} & \chgset & The bearer token the modules send each other. At least 32 characters; generate one with \val{openssl rand -hex 32}. \\
\opt{quarkus.flyway.migrate-at-start} & \val{true} & \chgkeep & Survey creates and upgrades its schema when it starts. \\
\opt{quarkus.datasource.jdbc.url} & \val{jdbc:postgresql://db:5432/survey} & \chgset & The database Survey connects to. A site points it at its own cluster. \\
\opt{quarkus.datasource.username} & \val{survey_user} & \chgmay & The role Survey connects as; it owns nothing and reads and writes the \val{survey} schema. \\
\opt{elicit.service.token} & \val{(none)} & \chgset & The bearer token the modules send each other. At least 32 characters; generate one with \val{openssl rand -hex 32}. \\
\opt{quarkus.flyway.migrate-at-start} & \val{true} & \chgkeep & Survey creates and upgrades its schema when it starts. \\
\opt{quarkus.datasource.jdbc.url} & \val{jdbc:postgresql://db:5432/survey} & \chgset & The database Survey connects to. A site points it at its own cluster. \\
\opt{quarkus.datasource.username} & \val{survey_user} & \chgmay & The role Survey connects as; it owns nothing and reads and writes the \val{survey} schema. \\
\opt{elicit.service.token} & \val{(none)} & \chgset & The bearer token the modules send each other. At least 32 characters; generate one with \val{openssl rand -hex 32}. \\
\opt{quarkus.flyway.migrate-at-start} & \val{true} & \chgkeep & Survey creates and upgrades its schema when it starts. \\
\opt{quarkus.datasource.jdbc.url} & \val{jdbc:postgresql://db:5432/survey} & \chgset & The database Survey connects to. A site points it at its own cluster. \\
\opt{quarkus.datasource.username} & \val{survey_user} & \chgmay & The role Survey connects as; it owns nothing and reads and writes the \val{survey} schema. \\
\opt{elicit.service.token} & \val{(none)} & \chgset & The bearer token the modules send each other. At least 32 characters; generate one with \val{openssl rand -hex 32}. \\
\opt{quarkus.flyway.migrate-at-start} & \val{true} & \chgkeep & Survey creates and upgrades its schema when it starts. \\
\opt{quarkus.datasource.jdbc.url} & \val{jdbc:postgresql://db:5432/survey} & \chgset & The database Survey connects to. A site points it at its own cluster. \\
\opt{quarkus.datasource.username} & \val{survey_user} & \chgmay & The role Survey connects as; it owns nothing and reads and writes the \val{survey} schema. \\
\opt{elicit.service.token} & \val{(none)} & \chgset & The bearer token the modules send each other. At least 32 characters; generate one with \val{openssl rand -hex 32}. \\
\opt{quarkus.flyway.migrate-at-start} & \val{true} & \chgkeep & Survey creates and upgrades its schema when it starts. \\
\end{settings}

\begin{plaintable}{tab:ports}{Published host ports in the shipped compose file}{@{\hspace{4pt}}>{\raggedright\arraybackslash}p{0.13\linewidth} >{\raggedright\arraybackslash}p{0.28\linewidth} >{\raggedright\arraybackslash}p{0.49\linewidth}@{\hspace{4pt}}}{\thead{Host port} & \thead{Service} & \thead{Reached by}}
\val{8080} & Survey & respondents, from a browser \\
\val{8081} & Admin & administrators, from a browser \\
\val{8082} & FHHS & Survey and Admin; the Family History definition names it \\
\val{8083} & Pedigree & FHHS only \\
\val{8084} & Author & authors, from a browser \\
\val{8085} & Author preview & authors, from a browser \\
\val{8180} & Keycloak & every browser and every module \\
\val{8025} & Mailpit & whoever reads the test mail \\
\val{16686} & Jaeger & whoever reads the traces \\
\val{5452} & PostgreSQL & the host, for \val{psql} \\
\end{plaintable}

\end{document}
